% Independent complete derivation; requires amsmath, amssymb, amsthm,
% amscd, hyperref, and proposition/lemma environments. Prefix FSC.
\section{A common sheaf for spherical, tropical, and full support coefficients}

\subsection{Sources and the exact base}
The human sources are Alain Connes and Caterina Consani,
\href{https://arxiv.org/abs/1502.05580}{\emph{Geometry of the Arithmetic
Site}}, original \nolinkurl{arithmeticsite_Adv_final1.tex}, specifically
Definition \texttt{site}, Theorem \texttt{structure2}, Remark
\texttt{caseno}, Theorem \texttt{thmspz}, and Proposition
\texttt{sheavesonspz0}; and their
\href{https://arxiv.org/abs/2606.06604}{\emph{On the Absolute Geometry
of Spec Z and the Fargues--Fontaine curve}}, original \nolinkurl{FF.tex},
Proposition \texttt{prop12}, the monoid-algebra adjunction in
Section~2.3, and the Frobenius formula at source lines 150--152.
The arithmetic-site source was read here at lines 598--704 and
900--1017, without a claim to exhaustive reading; FF was previously
read completely, lines 1--1764. The spherical functor $H$ is the
explicit coordinate functor of Connes--Consani's
\href{https://arxiv.org/abs/1502.05585}{\emph{Absolute algebra and
Segal's $\Gamma$-rings}}, original \nolinkurl{F1_corr.tex}, formula
\texttt{functsum} and Lemma \texttt{sssalg}. That source was previously
read completely, lines 1--1336.
The coefficient foundation is \emph{The Clankers}, \emph{Split Support
Geometry}, Version 11, Definition \texttt{def:lattice-split} and its
complete ideal and prime classification, in the
\href{https://github.com/KokunoYumeto/zeta-function-research-reader/blob/a2851f9eb352e7e4376bc629de86fa4ed9c1c434/workbenches/splitzero-tandem/foundations/split-support-geometry-v11/split_support_geometry_arithmetic_curve_v11.tex}{complete pinned original source}.
The constructions below give the common diagram for the companion
calculations ACF and ASL, with the required definitions and proofs
included here.

Let $L$ be a nontrivial bounded distributive lattice. It need not be
finite or Boolean. For a nonzero commutative ring $R$ retain
\begin{equation}\tag{FSC1}
 \begin{gathered}
 G_L(R)=\{(0,\lambda):\lambda\in L\}
                  \cup\{(a,1_L):a\in R\},\\
 (a,\lambda)+(b,\mu)=(a+b,\lambda\vee\mu),\qquad
 (a,\lambda)(b,\mu)=(ab,\lambda\wedge\mu),\\
 z_\lambda=(0,\lambda),\quad\tau=z_{0_L},\quad e=z_{1_L},
 \quad 1=(1_R,1_L),\quad q(a,\lambda)=a,\quad
 \sigma(a,\lambda)=\lambda.
 \end{gathered}
\end{equation}
Put $S=G_L(\mathbb Z)$, $X=\operatorname{Spec}S$,
$Y=\operatorname{Spec}\mathbb Z$, and
$A=\operatorname{Spec}_{\mathrm{lat}}L$. The additive semiring prime
spectrum has points
\begin{equation}\tag{FSC2}
 P_J=\{z_\lambda:\lambda\in J\}\quad(J\in A),\qquad
 Q_{\mathfrak p}=q^{-1}(\mathfrak p)\quad(\mathfrak p\in Y).
\end{equation}
For completeness, an ideal containing a top-supported element contains
$e$: add its product by $(-1,1_L)$ to itself if its amplitude is nonzero.
Then $ez_\lambda=z_\lambda$ puts the whole zero fibre in it, and its
amplitudes form a ring ideal. All other ideals are exactly the proper
lattice ideals in the zero fibre. For an amplitude ideal primality is
ring primality; for a support ideal it is meet primality. This proves
FSC2 directly. Their basic opens are
\begin{equation}\tag{FSC3}
 D_S(z_\lambda)=D_L(\lambda),\qquad
 D_S((n,1_L))=A\cup i(D_{\mathbb Z}(n)),\qquad
 A=D_S(e),\quad i(Y)=V_S(e).
\end{equation}
Here $i(\mathfrak p)=Q_{\mathfrak p}$ is the closed inclusion and
$j:A\hookrightarrow X$ the open inclusion. Every open meeting $i(Y)$
is uniquely $W_U=A\cup i(U)$ for a nonempty open $U\subseteq Y$;
the other opens are exactly the opens $V\subseteq A$. This follows
from FSC3 because every nonempty open of $Y$ contains its generic
point and all support points are generizations of every arithmetic
point. There is a continuous retraction
\begin{equation}\tag{FSC4}
 r(P_J)=r(Q_{(0)})=\eta,\qquad r(Q_{(p)})=(p),\qquad ri=\mathrm{id}.
\end{equation}
Indeed $r^{-1}D_{\mathbb Z}(n)=D_S((n,1_L))$ for $n\ne0$;
$r^{-1}\varnothing=\varnothing$. These constitute the basis test.
In particular $r^{-1}U=W_U$ for every nonempty open $U$.

\subsection{The support sheaf, its direct image, and all their sections}
For a lattice prime $J$ define
\begin{equation}\tag{FSC5}
 L_J=L[(L\setminus J)^{-1}],\qquad
 [\lambda]_J=[\mu]_J\iff
 \exists\nu\notin J:\lambda\wedge\nu=\mu\wedge\nu.
\end{equation}
Every invertible idempotent is the unit; hence this formula follows
from the localization equality, and every fraction is represented by
its numerator. It makes explicit the possible loss of a label in an
individual support stalk.

\begin{lemma}[Sections of the support sheaf]
The restriction $\mathcal O_A=j^{-1}\mathcal O_X$ of the actual
localization sheaf of $S$ satisfies
\begin{equation}\tag{FSC6}
 \mathcal O_A(D_L(\lambda))=\downarrow\lambda,\qquad
 a|_{D_L(\mu)}=a\wedge\mu\quad(\mu\le\lambda),\qquad
 (\mathcal O_A)_J=L_J.
\end{equation}
In particular $\Gamma(A,\mathcal O_A)=L$. For every open $V\subseteq A$
its sections are exactly
\begin{equation}\tag{FSC7}
 \left\{(a_\lambda)_{D(\lambda)\subseteq V}:
 a_\lambda\le\lambda,\quad
 a_\lambda\wedge\mu=a_\mu\text{ whenever }\mu\le\lambda
                         \text{ and }D(\lambda)\subseteq V\right\}.
\end{equation}
The two operations are coordinatewise join and meet.
\end{lemma}
\begin{proof}
First $S[e^{-1}]=L$ because $e^2=e$ makes $e=1$ and
$e(a,\lambda)=z_\lambda$. Then
$L[\lambda^{-1}]=\downarrow\lambda$ by $a\mapsto a\wedge\lambda$,
with unit $\lambda$. Indeed the equality of localized elements is
exactly equality after meeting with $\lambda$.
To verify the sheaf assertion, the prime-ideal separation theorem for
bounded distributive lattices gives an injective map
$\lambda\mapsto D_L(\lambda)$ preserving finite joins and meets, and
these basic opens are compact. Here is the needed argument.
If an ideal and filter are disjoint, an ideal maximal subject to
avoiding the filter exists by Zorn's lemma and is prime: if it contains
$a\wedge b$ but contains neither $a$ nor $b$, adjoining either meets
the filter, and distributivity makes the meet of those two filter
elements belong to the original ideal, a contradiction.
Apply this with the ideal $\downarrow\mu$ and filter
$\uparrow\lambda$ when $\lambda\nleq\mu$ to obtain separation.
For compactness, a cover $D(\lambda)\subseteq\bigcup_iD(\mu_i)$
without a finite subcover would leave the ideal generated by all
$\mu_i$ disjoint from $\uparrow\lambda$. The same prime would
contradict the cover.
For a finite cover $\lambda=\bigvee_i\lambda_i$, sections
$a_i\le\lambda_i$ agree on overlaps precisely when
$a_i\wedge\lambda_j=a_j\wedge\lambda_i$. Their unique gluing is
$\bigvee_i a_i$, by distributivity. Compactness reduces any cover of a
basic open to this case. Thus FSC6 defines the localization sheaf on
a basis and has the stalk FSC5. Restricting a section to every basis
open gives FSC7; conversely those compatible values glue and are
unique. This proves the formula for all opens, including $V=A$.
\end{proof}

Set $\mathcal D=j_*\mathcal O_A$. There is an actual unital restriction
map $d:\mathcal O_X\to\mathcal D$. Its sections and stalks are
\begin{equation}\tag{FSC8}
 \begin{array}{c|cc}
 &W_U=A\cup i(U)&V\subseteq A\\ \hline
 \mathcal D&L&\mathcal O_A(V)\\
 \mathcal O_X&G_L(\mathcal O_Y(U))&\mathcal O_A(V)
 \end{array}
 \qquad
 \mathcal D_{Q_{\mathfrak p}}=L,\quad\mathcal D_{P_J}=L_J.
\end{equation}
On $W_U$, $d$ is $\sigma$; on $V$, it is the identity.
To justify the remaining entry for $\mathcal O_X$, localizing at an
arithmetic prime gives $G_L(\mathbb Z_{\mathfrak p})$, including
$G_L(\mathbb Q)$ at $Q_{(0)}$: all denominators there are nonzero
integer amplitudes with top support, so labels never change. A section
over $W_U$ has one rational amplitude and one label at $Q_{(0)}$.
Generization from each $Q_{(p)}$ injects
$G_L(\mathbb Z_{(p)})\hookrightarrow G_L(\mathbb Q)$, forcing its
amplitude to lie in $\mathcal O_Y(U)$ and its label to be the same.
At all support points the section is determined by the germ at
$Q_{(0)}$, and equals the localization of that label. Conversely the
corresponding element of $G_L(\mathcal O_Y(U))$ gives these local
fractions. This proves FSC8. Every neighbourhood of an arithmetic
point contains all of $A$, so the direct-image stalk there is $L$;
at a support point it is the ordinary support stalk. Restriction
from $W_U$ to $V$ sends $(a,\lambda)$ to the section represented by
$\lambda$, whose component on $D(\mu)\subseteq V$ is
$\lambda\wedge\mu$.

Write $\underline L_X$ for the constant sheaf: its sections are the
locally constant maps into the discrete set $L$, with pointwise lattice
operations. There is a canonical unital sheaf map
\begin{equation}\tag{FSC9}
 \ell:\underline L_X\longrightarrow\mathcal D,
 \quad \ell_{Q_{\mathfrak p}}=\mathrm{id}_L,
 \quad \ell_{P_J}(\lambda)=[\lambda]_J.
\end{equation}
On $W_U$ a locally constant map is constant: every point of $A$ is a
generization of $Q_{(0)}$ and every other arithmetic point specializes
from $Q_{(0)}$, whereas a locally constant map into a discrete space
is constant along either relation. Thus its sections there are $L$.
On $V\subseteq A$ send a locally constant label function to the
section whose germ at $J$ is the localized label there; it is locally
represented by a single label and therefore is a section of $\mathcal
O_A$. This proves FSC9 on sections and under restrictions.

\subsection{Tropical coefficients, labels, and Frobenius zero}
For each rational prime $p$, put $H_p=\mathbb Z[1/p]$ and
$K_p=H_{p,\max}=H_p\cup\{-\infty\}$ with operations maximum and
ordinary addition. Its semiring zero is $\mathbf0=-\infty$, its
semiring unit is $\mathbf1=0\in H_p$. Define
\begin{equation}\tag{FSC10}
 K_{p,L}=K_p\otimes_{\mathbb B}L,\quad
 \kappa(\lambda)=\mathbf1\otimes\lambda,\quad
 \rho\left(\bigoplus_i h_i\otimes\lambda_i\right)
       =\bigvee_{h_i\ne-\infty}\lambda_i.
\end{equation}
The map $K_p\to\mathbb B$ sending every finite exponent to $1$ and
$-\infty$ to $0$ is a semiring homomorphism. Tensoring it proves that
$\rho$ is well defined and a homomorphism; the tensor relations prove
the same for $\kappa$, and $\rho\kappa=\mathrm{id}$.
Consequently $\kappa$ is injective without collapsing any label.
Positive Frobenius multiplies every finite $h_i$ by its integer
index $n$, while fixing $-\infty$. Exponent-zero Frobenius sends every
finite exponent to $0$ and fixes $-\infty$, so
\begin{equation}\tag{FSC11}
 \operatorname{Fr}_0=\kappa\rho,\qquad
 \rho\operatorname{Fr}_n=\rho\quad(n\ge0),\qquad
 \operatorname{Fr}_n\kappa=\kappa.
\end{equation}
These are semiring maps: scaling preserves maximum and exponent
addition, and the zero action preserves which coefficients are
absorbing zero. The last assertions also follow directly from the
tensor formulas.

The source geometric morphism $\Theta:Y\to\widehat{\mathbb N_0^\times}$
has $\Theta(\eta)=\{0\}$ and $\Theta(p)=H_p$.
Let $\mathcal T_L=\Theta^*(\mathcal O^0\otimes_{\mathbb B}\underline L)$.
The source Proposition \texttt{sheavesonspz0}, including the zero
Frobenius convention in Remark \texttt{caseno}, gives
\begin{equation}\tag{FSC12}
 (\mathcal T_L)_\eta=L,\qquad(\mathcal T_L)_p=K_{p,L}.
\end{equation}
The tensor commutes with inverse images: its finite products,
coproducts, and coequalizers imposing the semiring relations are
preserved by those functors. More explicitly, for every nonempty open
$U\subseteq Y$ the complete section description is
\begin{equation}\tag{FSC13}
 \begin{gathered}
 \mathcal T_L(U)=\{(\lambda,(w_p)_{p\in U}):
       \lambda\in L, w_p\in K_{p,L},\ \rho(w_p)=\lambda,\\
       w_p=\kappa(\lambda)\text{ for all but finitely many primes }p\in U\}.
 \end{gathered}
\end{equation}
Operations are componentwise and restrictions discard removed prime
coordinates. To prove this rather than infer a sheaf from stalks,
use the source description: an original coefficient section is either
absorbing zero everywhere, or finite exponents with value $0$ outside
a finite set of primes. Each local tensor expression therefore has a
single generic label $\lambda$ and is $\kappa(\lambda)$ outside a
finite set. Every open of $Y$ is quasi-compact, since a single nonempty
member of a cover misses only finitely many points. A finite cover of
local expressions gives precisely FSC13. Conversely, for each prime
in its finite exceptional set, express $w_p$ as a finite tensor sum;
extend each finite coefficient to the source section with that value
at $p$ and exponent $0$ elsewhere. On the open excluding the other
exceptional primes this realizes $w_p$ and is $\kappa(\lambda)$
elsewhere. These sections agree on overlaps and glue uniquely. Zero
coefficients can be omitted, including the all-zero case. This proves
both directions and every stated restriction map.

Set $\mathcal E=r^{-1}\mathcal T_L$. It has maps
$\rho:\mathcal E\to\underline L_X$ and
$\kappa:\underline L_X\to\mathcal E$ with $\rho\kappa=\mathrm{id}$.
Its sections on the two types of open, and its stalks, are
\begin{equation}\tag{FSC14}
 \begin{gathered}
 \mathcal E(W_U)=\mathcal T_L(U),\qquad
 \mathcal E(V)=\operatorname{LC}(V,L)\quad(V\subseteq A),\\
 \mathcal E_{Q_{(p)}}=K_{p,L},\quad
 \mathcal E_{Q_{(0)}}=L,\quad\mathcal E_{P_J}=L.
 \end{gathered}
\end{equation}
Here $\operatorname{LC}$ denotes locally constant maps. For the first
identity the inverse-image adjunction gives a map in one direction
and restriction along $i$ a left inverse. Sections with equal
arithmetic restriction have equal germs at $Q_{(0)}$, hence at every
support generization, so the left inverse is also injective. For the
second, $r|_A$ factors through the point $\eta$; inverse image to that
point is the stalk $L$, and pulling it back to $A$ is the constant
sheaf. The stalk formula follows either way. Restriction from $W_U$
to $V$ remembers the constant label $\lambda$ in FSC13. Finally put
\begin{equation}\tag{FSC15}
 t=\ell\rho:\mathcal E\longrightarrow\mathcal D.
\end{equation}
This is the promised unital tropical-to-support map: at arithmetic
points it is $\rho$, at $P_J$ it is $L\to L_J$, and on all sections
it is fixed by every Frobenius, including zero.


\subsection{The actual fibre product and its constrained carrier}
Form the following fibre product of sheaves of unital semirings:
\begin{equation}\tag{FSC16}
 \mathcal C=\mathcal O_X\mathbin{\times}_{\mathcal D}\mathcal E,
 \qquad
 \mathcal C(W)=\{(s,w):d(s)=t(w)\text{ in }\mathcal D(W)\}.
\end{equation}
All operations are componentwise. Compatible pairs glue uniquely in
the two factor sheaves, and the equality in $\mathcal D$ can be
checked locally, so FSC16 is already a sheaf. The projections are
unital homomorphisms. For any other sheaf with maps to both factors
whose composites to $\mathcal D$ agree, their pair is the unique
map to FSC16. This proves the fibre-product universal property.

Its sections on every open and its stalks are as follows:
\begin{equation}\tag{FSC17}
 \begin{gathered}
 \mathcal C(W_U)=
 \{(a,\lambda,(w_p)):\ a\in\mathcal O_Y(U),\
            (\lambda,(w_p))\in\mathcal T_L(U),\
            a\ne0\Rightarrow\lambda=1_L\},\\
 \mathcal C(V)=\operatorname{LC}(V,L)\quad(V\subseteq A).
 \end{gathered}
\end{equation}
In the second line a label function $f$ corresponds to the pair
$(\ell f,f)$, not only to its localized value $\ell f$.
Restrictions among the $W_U$ localize the amplitude and remove prime
coordinates. Restriction to $V\subseteq A$ forgets the amplitude and
prime coordinates and keeps the constant function $\lambda$.
On the arithmetic finite stalk, abbreviate $R_p=\mathbb Z_{(p)}$.
Then
\begin{equation}\tag{FSC18}
 \begin{gathered}
 \mathcal C_{Q_{(p)}}
  \simeq\{(a,w)\in R_p\times K_{p,L}:a\ne0\Rightarrow\rho(w)=1_L\},\\
 \mathcal C_{Q_{(0)}}\simeq G_L(\mathbb Q),\qquad
 \mathcal C_{P_J}\simeq L.
 \end{gathered}
\end{equation}
The finite-stalk pair $(a,w)$ corresponds to
$((a,\rho(w)),w)$ in the original fibre product. Its sum and product
are $(a+b,w\oplus v)$ and $(ab,wv)$, with their exact original
amplitudes and tropical coefficients. The support stalk $\lambda$
corresponds to $([\lambda]_J,\lambda)$.
To prove these assertions, substitute FSC8 and FSC14 into FSC16.
For stalks, every finite set of section representatives can be put on
one common neighbourhood, and an equality of their germs holds on
a smaller common neighbourhood. Thus a stalk commutes with this finite
fibre product. At a support prime the two maps to $L_J$ are the identity
on $L_J$ and the localization $L\to L_J$, so their fibre product is
$L$, by precisely the pair displayed above. At an arithmetic point
the two maps are $\sigma$ and $\rho$, which give the first two lines.
This proves all sections and stalks rather than identifying the three
input sheaves.

Write $\mathcal R=i_*\mathcal O_Y$. On $W_U$ it is
$\mathcal O_Y(U)$ and on every $V\subseteq A$ it is the zero ring
(the singleton ring with $0=1$). The amplitude of an actual local
fraction gives a unital map $q:\mathcal O_X\to\mathcal R$.
There is an injective unital comparison
\begin{equation}\tag{FSC19}
 \mathcal C\hookrightarrow\mathcal R\times\mathcal E,
              \qquad (s,w)\longmapsto(q(s),w).
\end{equation}
At finite arithmetic stalks its image is exactly the constrained
subset in FSC18; at the generic arithmetic stalk it is
$G_L(\mathbb Q)\subset\mathbb Q\times L$; at support stalks it is
all of $\{0\}\times L$. Injectivity follows directly from these
stalk descriptions: on arithmetic stalks the amplitude and
$\rho(w)$ recover $s$, while on support stalks $s$ is the localization
of $w$. A sheaf map injective on stalks is injective on sections,
since equality of germs gives local equality and then equality.
The additional elements in the unrestricted finite-stalk ambient
product are exactly
\begin{equation}\tag{FSC20}
 \{(a,w):a\ne0,\ \rho(w)\ne1_L\}.
\end{equation}
They are not present in $\mathcal C$. At the generic stalk the
corresponding extra pairs are $a\ne0$ with support below $1_L$.
Thus the correspondence never silently enlarges the original
nonzero-amplitude support constraint.

\subsection{Original coefficients, the tropical ideal, and their maps}
Let $B_L$ be the coefficient sheaf on $Y$ with
$B_L(U)=G_L(\mathcal O_Y(U))$ for nonempty $U$. A compatible cover has
one support label because all nonempty intersections contain $\eta$;
the ring amplitudes glue and retain top support when nonzero. This
proves the sheaf assertion. Preserve the actual inverse image under
the separate name
\[
 \mathcal P=r^{-1}B_L.
\]
Its sections on $W_U$ are $G_L(\mathcal O_Y(U))$, while its sections
on $V\subseteq A$ are $\operatorname{LC}(V,G_L(\mathbb Q))$.
Its arithmetic stalks are $G_L(R_p)$ and $G_L(\mathbb Q)$, and its
stalk at \emph{every} support point is also $G_L(\mathbb Q)$.
These follow from the inverse-image stalk formula and the factorization
of $r|_A$ through $\eta$, whose $B_L$-stalk is $G_L(\mathbb Q)$.
Thus this original pullback still retains rational amplitudes on the
support open. It is not the fixed sheaf constructed next.

Define the separate sheaf
\[
 \mathcal B=\mathcal O_X\mathbin{\times}_{\mathcal D}\underline L_X.
\]
Its sections on $W_U$ are $G_L(\mathcal O_Y(U))$, represented by
$((a,\lambda),\lambda)$; its sections on $V\subseteq A$ are
$\operatorname{LC}(V,L)$, represented by $(\ell f,f)$.
Its stalks are $G_L(R_p)$, $G_L(\mathbb Q)$, and $L$ at support
points. These formulas are direct substitutions from FSC8--FSC9,
with the same sheaf and finite-stalk argument as for FSC16--FSC18.
There is an exact sequence of comparison maps (no additive exactness
terminology is imposed)
\begin{equation}\tag{FSC33}
 \mathcal P\xrightarrow{\alpha}\mathcal B
                       \xrightarrow{\Phi}\mathcal O_X.
\end{equation}
Here $\Phi$ is the first projection. On arithmetic stalks both maps
are the identity under the displayed identifications; at $P_J$ they
are exactly
\begin{equation}\tag{FSC34}
 G_L(\mathbb Q)\xrightarrow{\sigma}L
                   \longrightarrow L_J,\qquad
 (a,\lambda)\longmapsto\lambda\longmapsto[\lambda]_J.
\end{equation}
To construct them on sections, let $\alpha$ be the identity on $W_U$
and the coordinatewise support map on locally constant sections
$V\to G_L(\mathbb Q)$. Restriction of an arithmetic-open section to
$V$ is its constant rational-amplitude germ in $\mathcal P$, and
its constant support label in $\mathcal B$; hence these formulas
commute with that restriction as well as restrictions within each
kind of open. Operations are preserved by $\sigma$. This proves
$\alpha$ as an actual unital sheaf map and proves FSC34.
Both maps are epimorphisms because all their stalk maps are surjective.
The equality relation of $\alpha$ at support points is equality of
support labels, with amplitudes arbitrary within the original
constrained carrier; at arithmetic points it is equality of the whole
element. The equality relation of $\Phi$ at support points is FSC5,
and at arithmetic points again equality. Their composite is the
previous full coefficient-to-localization map, with the support-only
loss of amplitude occurring at the first arrow and lattice localization
at the second. All three sheaves remain explicit.

There is an injective unital map and a unital retraction
\begin{equation}\tag{FSC21}
 \begin{gathered}
 \nu:\mathcal B\hookrightarrow\mathcal C,\qquad
 (a,\lambda)\longmapsto(a,\kappa(\lambda)),\\
 \beta:\mathcal C\longrightarrow\mathcal B,\qquad
 (a,w)\longmapsto(a,\rho(w)),\qquad \beta\nu=\mathrm{id}.
 \end{gathered}
\end{equation}
The displayed formulas apply at arithmetic stalks and to the
corresponding prime coordinates of sections; at support points
both maps are the identity of $L$. They respect every restriction
formula already proved. Their operations follow because $\kappa$
and $\rho$ are homomorphisms and $\rho\kappa=\mathrm{id}$.
The composite $\mathcal B\to\mathcal C\to\mathcal O_X$ is the
isomorphism at arithmetic stalks and the localization $L\to L_J$
at support stalks. These are the exact places where the fixed
coefficient comparison and the actual localization sheaf differ.

Frobenius on the common object is
\begin{equation}\tag{FSC22}
 F_n(s,w)=(s,\operatorname{Fr}_n w)\quad(n\ge0),\qquad
 F_0=\nu\beta,
 \qquad \operatorname{Fix}(F_0)=\nu(\mathcal B).
\end{equation}
It is a sheaf endomorphism because $\rho\operatorname{Fr}_n=\rho$.
Every amplitude is retained unchanged; no claim is made that ordinary
ring powering is an additive endomorphism. The last two assertions
follow from FSC11 and FSC21. They also recover the original constrained
$G_L(R_p)$ exactly at every arithmetic stalk. As sheaves of semirings,
$F_0$ is an actual idempotent with the indicated image, not an assumed
identification of the spherical, tropical and arithmetic structures.

There is a second exact map
\begin{equation}\tag{FSC23}
 \delta:\mathcal E\longrightarrow\mathcal C,\qquad
 \delta(w)=(z_{\rho(w)},w)	ext{ at arithmetic stalks},\qquad
 \delta(\lambda)=\lambda\text{ at support stalks}.
\end{equation}
On $W_U$ it sets the amplitude equal to zero and retains all tropical
coordinates. On $V$ it is the identity on locally constant labels.
These descriptions commute with restrictions and give an injective
map preserving zero, addition, and multiplication. It is nonunital:
its unit image has arithmetic amplitude zero. More precisely let
$b=\delta(1_{\mathcal E})$, the global section whose pair in FSC16 is
$(e,1_{\mathcal E})$. Then
\begin{equation}\tag{FSC24}
 b^2=b,\qquad
 b\mathcal C=\delta(\mathcal E)
          =\ker(\mathcal C\longrightarrow\mathcal R),\qquad
 \delta:\mathcal E\xrightarrow{\sim}b\mathcal C.
\end{equation}
The last map is unital when the principal ideal is given its own
identity $b$. Indeed multiplication by $b$ sends $(a,w)$ to $(0,w)$,
which proves both equalities and the asserted identity on the ideal;
at support points $b=1$ and the same formula holds in the zero-amplitude
ambient description. All inverse maps are given by the tropical
projection. There cannot be a unital sheaf map
$\mathcal E\to\mathcal C$: addition is idempotent in $\mathcal E$,
whereas $1_{\mathcal C}+1_{\mathcal C}\ne1_{\mathcal C}$ at the
arithmetic generic stalk, since its amplitudes are $2$ and $1$ in
$\mathbb Q$. Thus FSC23--FSC24 are the exact replacement for the
invalid unital inclusion, while retaining the entire ideal.

The amplitude map is the ring reflection of this common sheaf:
\begin{equation}\tag{FSC25}
 \mathcal C/\!\sim_{\delta\mathcal E}\ \simeq\mathcal R,
 \qquad x\sim_{\delta\mathcal E}y
 \iff \text{locally }x+u=y+v\text{ for }u,v\in\delta\mathcal E.
\end{equation}
It is surjective on sections of $W_U$, since $a$ has a lift
$(a,\kappa(1_L))$, and on sections of support opens, whose target is
one element. If two arithmetic-open sections $(a,w)$ and $(a,w')$
have the same amplitude, then
\[
 (a,w)+\delta(w')=(a,w')+\delta(w)
\]
in the full pair presentation: supports join on both sides and the
tropical sums agree. At support opens the whole semiring is the ideal.
Conversely adding ideal sections does not change amplitude. This
proves the exact congruence on all opens, hence its sheaf version.
Any unital map to a sheaf of rings annihilates this ideal because
its elements are additively idempotent; it therefore factors uniquely
through FSC25. This proves the universal property, without any
unsupported stable-spectrum consequence.

\subsection{The spherical map into the entire common object}
Let $\mathcal F_L=r^{-1}\Theta^{-1}\mathbf F_1[T]$, with
$\mathcal F_\eta=\mathbf F_1$ and
$\mathcal F_p=\mathbf F_1[T^{H_p^+}]$ as in FF Proposition
\texttt{prop12}. On the source topos the pointed multiplicative map
\begin{equation}\tag{FSC26}
 \mathbf F_1[T]\longrightarrow H(\mathbb Z_{\max}\otimes_{\mathbb B}L),
 \qquad T^n\longmapsto n\otimes1_L,\qquad
 0\longmapsto\mathbf0
\end{equation}
preserves the unit: $T^0$ goes to exponent $0$, the tropical unit,
not to the absorbing $-\infty$. It preserves multiplication because
$n+m$ is the tropical product, and it commutes with every source
Frobenius including zero. To see its extension to spherical algebras
explicitly, a monoid-algebra element in degree $k_+$ has a single
monomial and one nonbasepoint coordinate. Put its image in that
coordinate and the additive zero elsewhere. Pointed maps move or
discard that coordinate, agreeing with the fibre-sum formula of $H$;
products agree by the proved multiplicative formula. Conversely
coordinate inclusions force this construction. Thus FSC26 is an
actual spherical-algebra morphism. Inverse image gives
\begin{equation}\tag{FSC27}
 u:\mathcal F_L\longrightarrow H\mathcal E,
 \qquad u_{Q_{(p)}}(T^h)=h\otimes1_L,
 \quad u_{P_J}:\mathbf F_1\to HL\text{ the unit map}.
\end{equation}
The same unit formula holds at $Q_{(0)}$.

There is also the augmentation
$v:\mathcal F_L\to H\mathcal O_X$, with $T^h\mapsto1$ and the
absorbing zero mapped to the target additive zero. It exists by
pulling back the equivariant source augmentation $T^n\mapsto1$.
It is unique: any such map is the unit map at $Q_{(0)}$, whose
source is $\mathbf F_1$; the injection
$G_L(R_p)\to G_L(\mathbb Q)$ then forces every monomial's image at
$Q_{(p)}$ to be $1$. The source generization has $T^h\mapsto1$,
because it commutes with $\operatorname{Fr}_0$, which is the identity
on $\mathbf F_1$ and sends $T^h$ to $T^0$. At support primes the source
is also $\mathbf F_1$, so the unit map is unique there. Equality on
stalks proves equality of sheaf maps.

The composites $Hd\circ v$ and $Ht\circ u$ are equal. On every
finite stalk both send each monomial to $1_L$ and the absorbing zero
to $0_L$; at support points both give the unit map to $HL_J$.
The monoid-algebra coordinate argument makes this equality hold in
every finite pointed degree. Since $H$ takes the semiring fibre product
to the spherical fibre product degree by degree, there is a unique map
\begin{equation}\tag{FSC28}
 \begin{gathered}
 \gamma:\mathcal F_L\longrightarrow H\mathcal C,\\
 \gamma_{Q_{(p)}}(T^h)=(1,h\otimes1_L),\qquad
 \gamma_{Q_{(p)}}(0)=(0,\mathbf0),\\
 \gamma_{Q_{(0)}}(1)=(1,1_L),\quad
 \gamma_{Q_{(0)}}(0)=(0,0_L),\qquad
 \gamma_{P_J}:\mathbf F_1\to HL\text{ the unit map}.
 \end{gathered}
\end{equation}
Here $(1,h\otimes1_L)$ denotes exactly the constrained pair of FSC18.
This map is injective. At a finite stalk, distinct $h$ stay distinct:
a lattice prime exists by the separation proof above, and its Boolean
character gives an evaluation $K_{p,L}\to K_p$ taking
$h\otimes1_L$ to $h$. The absorbing zero has a different arithmetic
amplitude. At the generic and support stalks the two values are
also distinct, since $L$ is nontrivial. In each finite pointed degree
the coordinate description has at most one nonzero entry, preserving
its coordinate and value. Thus it is injective in every degree and
every stalk, hence as a sheaf map. This proves a faithful spherical
comparison retaining positive exponents, even though its projection
to the ordinary split structure sheaf is the forced augmentation.
Finally $F_n\gamma=\gamma\operatorname{Fr}_n$ for all $n\ge0$ by
FSC22 and the displayed monomial formulas; no ordinary amplitude
Frobenius has been introduced.

The complete diagram is
\begin{equation}\tag{FSC29}
\begin{CD}
 \mathcal F_L @>{\gamma}>> H\mathcal C @>>> H\mathcal E\\
 @V{v}VV @VVV @VV{Ht}V\\
 H\mathcal O_X @= H\mathcal O_X @>{Hd}>> H\mathcal D.
\end{CD}
\end{equation}
The right square is a fibre product, all arrows are actual unital
spherical sheaf maps, and the upper-left map is injective. The ideal
map $\delta$ of FSC23 is deliberately outside this diagram of unital
maps; it becomes a unital isomorphism only onto the ideal with identity
$b$ in FSC24. Exact proof locators for the topological base and bottom
row are FSC2--FSC9; for the right-hand maps FSC10--FSC16; and for the
spherical monomial maps FSC26--FSC28.


\subsection{The precise obstruction to a global section and its replacement}
The projection $\pi:\mathcal C\to\mathcal O_X$ is an epimorphism of
sheaves, but for the full lattice its arithmetic section need not
extend over the support open. Its exact stalk fibres are
\begin{equation}\tag{FSC30}
 \begin{gathered}
 \pi^{-1}_{Q_{(p)}}(a,\lambda)
       =\{w\in K_{p,L}:\rho(w)=\lambda\},\qquad
 \pi_{Q_{(0)}}=\mathrm{id}_{G_L(\mathbb Q)},\\
 \pi^{-1}_{P_J}([\mu]_J)
       =\{\lambda\in L:\exists\nu\notin J,                          \lambda\wedge\nu=\mu\wedge\nu\}.
 \end{gathered}
\end{equation}
The finite fibres are nonempty, containing $\kappa(\lambda)$;
the support fibres are nonempty because localization is represented
by a numerator. Thus it is surjective on stalks. To see directly that
this implies the sheaf epimorphism assertion, any section of the target
has a preimage germ at each point; choose local representative sections
of the domain and shrink so their images equal the given target
section. These local lifts prove local surjectivity, the epimorphism
condition for sheaves of sets with algebraic operations.

A unital sheaf section of $\pi$ exists if and only if $L$ is the
two-element lattice. First suppose one exists. Compose it with $\beta$
to obtain a section of $\Phi:\mathcal B\to\mathcal O_X$ from FSC21.
At $Q_{(0)}$, $\Phi$ is the identity, so this section is forced to be
the identity there. Compatibility with generization to $P_J$ forces
its support-stalk component to satisfy
\begin{equation}\tag{FSC31}
 s_J([\lambda]_J)=\lambda\qquad\text{for every }\lambda\in L.
\end{equation}
Indeed apply the compatibility square to $z_\lambda$ at $Q_{(0)}$:
in $\mathcal O_X$ it generizes to $[\lambda]_J$, whereas in
$\mathcal B$ it generizes to the unlocalized label $\lambda$.
If $0_L<\lambda<1_L$, prime-ideal separation supplies a prime $J$
avoiding $\lambda$. In $L_J$, $[\lambda]_J=[1_L]_J$, which contradicts
FSC31. Thus no intermediate element exists and $L=\{0_L,1_L\}$.
Conversely for this two-element lattice its sole prime is $\{0_L\}$,
and localization there is the identity. Therefore $\Phi$ is an
isomorphism at every stalk, hence a sheaf isomorphism: local inverses
exist and glue uniquely by injectivity. The map $\nu\Phi^{-1}$ then
is a unital section of $\pi$. This proves both directions.

The obstruction therefore defines the exact replacement rather than
removing the support geometry: $\mathcal B$ is the coefficient sheaf
which retains the unlocalized support labels, embeds by $\nu$ into
$\mathcal C$, and maps onto the actual $\mathcal O_X$ by precisely
FSC31's localization quotient. On the arithmetic closed stratum,
\begin{equation}\tag{FSC32}
 i^{-1}\mathcal C=B_L\mathbin{\times}_{\underline L_Y}\mathcal T_L,
 \qquad B_L\hookrightarrow i^{-1}\mathcal C,
 \quad (a,\lambda)\longmapsto(a,\kappa\lambda).
\end{equation}
To verify this, inverse image commutes with finite products and
equalizers, or one may substitute the already proved sections and
stalks and use the same finite-representative germ argument as in
FSC18. Here $i^{-1}\mathcal O_X=B_L$, $i^{-1}\mathcal D=\underline L_Y$,
and $i^{-1}\mathcal E=\mathcal T_L$. Thus the original constrained
coefficient sheaf survives exactly, while its globally extendable
replacement and the actual localization map remain visible.

For the three-element chain $0<a<1$, the support primes are
$J_0=\{0\}$ and $J_a=\{0,a\}$. At $J_0$, localization sends $a$ to
$1$ and gives $L_{J_0}=\{0,1\}$; at $J_a$, only $1$ is inverted and
$L_{J_a}=L$. Accordingly $\mathcal C$ has the three-element stalk
$L$ at both support points, whereas $\mathcal O_X$ has respectively
two and three elements. The fibre over $1\in L_{J_0}$ in FSC30 is
exactly $\{a,1\}$. This proves the obstruction concretely without
replacing the chain lattice by independent Boolean coordinates.
At a finite arithmetic stalk, $w=\mathbf1\otimes a$ has
$\rho(w)=a$. The pair $(1,w)$ therefore exists in the unrestricted
ambient product but is excluded by FSC20, whereas $(0,w)$ exists in
$\mathcal C$. The original three zero labels appear exactly as
$(0,\kappa0),(0,\kappa a),(0,\kappa1)$ in its fixed coefficient
subobject. These examples describe the map, its fibre, and the
excluded amplitudes using the same original lattice throughout.
